% !PRINTED 264
\weberpage{264}
%<<CONT
haben, dass aber in der Kette der Haupt-Unterdeterminanten
$R'_k$ der Form $\varphi'$ keine verschwindenden Glieder vorkommen, und
dass endlich einem nicht verschwindenden $R_k$ ein $R'_k$ von dem\-selben
Vorzeichen entspricht. Dann können die Zahlen $\pi$, $\nu$ für
$\varphi$ und für $\varphi'$ sowohl aus der Reihe (4), als auch aus der Reihe
(5) ermittelt werden, und die Anzahl der \gesperrt{Zeichenwechsel}, die
in beiden gleich ist, giebt die Anzahl $\nu$ der negativen Quadrate.

Um nun den Nachweis zu führen, dass der Coëfficient $\varepsilon$ in
der angegebenen Weise bestimmt werden kann, nehmen wir an,
es sei $\varphi$ irgendwie in eine Summe von Quadraten verwandelt
\[
\varphi = \lambda_1 y_1^2 + \lambda_2 y_2^2 + \cdots + \lambda_m y_m^2,
\]
und $\psi$, in den Variablen $y_1, y_2 \ldots y_m$ dargestellt, habe den
Ausdruck
\[
\psi = \Sigma \beta_{i,k} y_i y_k.
\]
Die Zahlen $\pi$, $\nu$ für die Function $\varphi'$ werden dann aus der
Kette der Haupt-Unterdeterminanten von
\[
\begin{vmatrix}
\lambda_1 + \varepsilon\beta_{1,1}, & \varepsilon\beta_{1,2} \ldots & \varepsilon\beta_{1,m} \\
\varepsilon\beta_{2,1}, & \lambda_2 + \varepsilon\beta_{2,2} \ldots & \varepsilon\beta_{2,m} \\
\cdot & \cdot & \cdot \\
\varepsilon\beta_{m,1}, & \varepsilon\beta_{m,2} \ldots & \lambda_m + \varepsilon\beta_{m,m}
\end{vmatrix}
\]
nach dem Satze V bestimmt.

Nun kann man aber $\varepsilon$ so klein annehmen, dass diese Haupt-
Unterdeterminanten \gesperrt{dem Zeichen nach} übereinstimmen mit
\[
\lambda_1\lambda_2\lambda_3 \ldots \lambda_m, \quad \lambda_1\lambda_2\lambda_3 \ldots \lambda_{m-1}, \quad \ldots, \quad \lambda_1\lambda_2, \quad \lambda_1, \quad 1,
\]
und dann ist die Anzahl der positiven und negativen Quadrate
von $\varphi'$ gleich der Anzahl der positiven und negativen unter den
Coëfficienten $\lambda$, von denen keiner verschwindet, d.\,h. die Zahlen
$\pi$ und $\nu$ sind für $\varphi$ und $\varphi'$ dieselben.

Sind nun die Coëfficienten von $\psi$, in den ursprünglichen
Variablen ausgedrückt, $b_{i,k}$, also
\[
\psi = \Sigma b_{i,k} x_i x_k,
\]
so ist eine der Haupt-Unterdeterminanten von $\varphi'$
\[
R'_k = \begin{vmatrix}
a_{1,1} + \varepsilon b_{1,1} & \ldots & a_{1,k} + \varepsilon b_{1,k} \\
\cdot & \cdot & \cdot \\
a_{k,1} + \varepsilon b_{k,1} & \ldots & a_{k,k} + \varepsilon b_{k,k}
\end{vmatrix},
\]
und ist also eine ganze rationale Function $k^{\mathrm{ten}}$ Grades von $\varepsilon$,
\[
R'_k = R_k + \varepsilon M_1 + \varepsilon^2 M_2 + \ldots \varepsilon^k M_k,
\]
%CONT>>
